\section{二维一侧：WZW 模型与共形块}

\subsection{三维与二维的接口}

上一讲末尾，三维理论的 Hilbert 空间被写成了 $\cM_J$ 上全纯截面的空间。
这一讲要说明：同一个空间，也是二维共形场论的“共形块”空间。
两个理论在\textbf{态空间}层面是同一个东西；这不是说它们作为量子场论等价，
而是说同一个向量空间有两种描述。

二维一侧的好处是：那里的语言（表示、共形权、中心荷、融合规则）是纯代数的，
算起来比三维路径积分容易得多。所以这一讲的任务是把二维这套词汇搭起来，
让第七讲的每一次手术计算都有代数依据。读者如果只想记住一句话，那就是：
三维理论“知道”纽结，是因为它的态空间等于共形块空间，而共形块的融合规则就是纽结的拆接关系。

\subsection{WZW 作用量与 Wess--Zumino 项}

二维一侧的基本对象是 Wess--Zumino--Witten 模型：场 $g:\Sigma\to G$ 取值在紧致群中，
作用量由动能项与一个“拓扑”项相加：
\[
S_{\mathrm{WZW}}[g]
=\frac{k}{8\pi}\int_\Sigma \Tr\left(g^{-1}\partial g\;g^{-1}\bar\partial g\right)\dd^2z
+\frac{k}{12\pi}\int_B \Tr\left((g^{-1}\dd g)^3\right),
\]
其中 $\Sigma$ 是闭黎曼面，$B$ 是它的一个三维填充（$\partial B=\Sigma$）。
第一项是普通的 $\sigma$ 模型动能；第二项叫 Wess--Zumino 项。

它为什么必须写成三维积分？因为被积的 $\Tr(g^{-1}\dd g)^3$ 是一个\textbf{三形式}，
在二维流形 $\Sigma$ 上根本无从积分。用三维填充 $B$ 是它唯一的存在方式。
代价是 $B$ 不唯一：若 $B'$ 是另一个填充，$B\cup(-B')$ 是闭三维流形，
于是
\[
S_{\mathrm{WZW}}[g]\Big|_{B'}-S_{\mathrm{WZW}}[g]\Big|_{B}
=\frac{k}{12\pi}\int_{B\cup(-B')}\Tr\left((g^{-1}\dd g)^3\right)
=\frac{k}{12\pi}\cdot 24\pi^2\,n=2\pi k\,n,
\]
其中 $n$ 是映射 $g$ 的绕数。这一步用了 $\pi_3(G)\cong\ZZ$（紧致单连通单群）：
三形式的积分除以 $24\pi^2$ 给整数。因此作用量本身依赖填充，
但路径积分的权重 $\exp(\ii S)$ 良好定义，当且仅当
\[
k\in\ZZ .
\]
这正是第二讲里 $k$ 的整性来源，只是换了一件外衣：那边是大规范变换使
$\Delta L=2\pi k\,m$，这边是换填充使 $\Delta S=2\pi k\,n$；
两个整数 $m,n$ 都是同一个 $\pi_3(G)\cong\ZZ$ 的绕数。所以两处“$k$ 必须整”不是两个要求，
而是同一个要求。

\subsection{仿射李代数}

WZW 模型的手征流是
\[
J(z)=k\,\partial g\,g^{-1}\qquad(\text{反全纯流 } \bar J=k\,\bar\partial g\,g^{-1}),
\]
写成 Lie 代数分量 $J^a(z)$，它满足算子积展开
\begin{equation}
J^a(z)J^b(w)\sim\frac{k\,\delta^{ab}}{(z-w)^2}+\frac{\ii f^{abc}J^c(w)}{z-w}. \tag{6.1}
\end{equation}
把它按模式展开 $J^a(z)=\sum_{n\in\ZZ}J^a_n z^{-n-1}$，就得到无穷维代数
\[
[J^a_m,J^b_n]=\ii f^{abc}J^c_{m+n}+k\,m\,\delta^{ab}\,\delta_{m+n,0}.
\]
第一项是 Lie 代数的模式版本，第二项是\textbf{中心扩张}：
它是数（与 $J$ 无关），只在 $m+n=0$ 时出现，系数就是 level $k$。
这个代数记作 $\hat{\mathfrak g}_k$，叫仿射 Kac--Moody 代数。
没有第二项时，代数只是 $\mathfrak g$ 值函数的环路代数，没有最高权表示论可言；
中心扩张是整个二维一侧能“量子化”的关键。

\subsection{表示论、Weyl 群与 level 截断}

仿射代数的表示用最高权来分类。取一个有限维 $\mathfrak g$ 的不可约表示 $R_\lambda$
（最高权 $\lambda$），要求
\[
J^a_0|\lambda\rangle=t^a_\lambda|\lambda\rangle,
\qquad J^a_n|\lambda\rangle=0\quad(n>0).
\]
由 $J^a_{-n}$（$n>0$）作用生成的模叫 Verma 模。Verma 模一般可约：
某些组合会被零向量（null vector）打成零。把这些零向量商掉，得到不可约模
$\hat V_{k,\lambda}$。

何时 $\hat V_{k,\lambda}$ 是\textbf{酉}的（即内积正定）？答案是\textbf{可积}条件：
\[
(\lambda,\theta)\le k
\qquad\text{对所有正根 }\theta ,
\]
等价地 $(\lambda,\theta_{\max})\le k$，其中 $\theta_{\max}$ 是最大根。
满足条件的 $\lambda$ 叫 level $k$ 的可积最高权。几何上，这些 $\lambda$ 落在
仿射 Weyl 群
\[
W_a=W\ltimes kQ^\vee
\]
作用下的\textbf{基本 alcove} 内；Weyl 群并不直接作用在 $\lambda$ 上，
而是作用在平移后的 $\lambda+\rho$ 上（$\lambda\mapsto w(\lambda+\rho)-\rho$）。
\textbf{这正是 $\rho$ 与 $k+h^\vee$ 反复出现的原因}：所有公式都活在“平移过的”坐标里，
平移量是 $\rho$，而 $k$ 被 $h^\vee$ 补成 $k+h^\vee$ 才与 Weyl 群相容。

对 $SU(2)$，条件化简为 $2j\le k$，于是可积表示是
\[
j=0,\ \frac12,\ 1,\ \dots,\ \frac k2,
\]
一共 $k+1$ 个。对 $SU(N)$，可积表示的个数是
$\binom{k+N-1}{N-1}$（$SU(2)$ 时回到 $k+1$）。
这个“$k+1$”后面会反复出现：它是环面 Hilbert 空间的维数，是 $S$ 矩阵的阶数，
也是 Verlinde 公式里求和项的个数。

\subsection{Sugawara 构造与中心荷}

仿射代数上可以构造一个 Virasoro 代数，办法是把流“平方”：
\begin{equation}
T(z)=\frac{1}{2(k+h^\vee)}\sum_a :J^a(z)J^a(z):, \tag{6.2}
\end{equation}
其中 $:\ :$ 是正规序（把产生算子放左边），$h^\vee$ 是对偶 Coxeter 数。
这个构造叫 Sugawara 构造。它的合理性来自一个直接但繁琐的 Wick 收缩计算：
把 (6.1) 代进 $T(z)T(w)$，只保留四次极点，会发现
\begin{itemize}[leftmargin=*]
  \item 两次双收缩贡献的量级是 $k^2\dim G$；
  \item 一次双收缩加一次单收缩贡献的量级是 $k\,h^\vee\dim G$。
\end{itemize}
若只用双收缩，得到的“中心荷”是 $k^2\dim G/(2(k+h^\vee)^2)$，
它不满足 Virasoro 代数（$T$ 与自身的 OPE 不能闭合）。
补上单收缩项之后两项合并，恰好给出
\begin{equation}
c=\frac{k\dim G}{k+h^\vee}. \tag{6.3}
\end{equation}
分母里的 $k+h^\vee$ 不是凑出来的，它来自上述两类收缩的比值；
换句话说，$h^\vee$ 是流代数的“背景荷”，没有它 Virasoro 代数不闭合。

对 $SU(2)$，$c=3k/(k+2)$：$k=1$ 给 $c=1$，$k\to\infty$ 给 $c\to3$，
即三个自由玻色子——这与“大 $k$ 时理论趋于自由”的图像一致。
一般地 $k\to\infty$ 时 $c\to\dim G$。

表示 $R_\lambda$ 的共形权同样带 $k+h^\vee$：
\begin{equation}
h_\lambda=\frac{(\lambda,\lambda+2\rho)}{2(k+h^\vee)}. \tag{6.4}
\end{equation}
对 $SU(N)$ 的定义表示，$\lambda=\omega_1$，$(\omega_1,\omega_1)=(N-1)/N$、
$2(\omega_1,\rho)=N-1$，代入得
\[
h=\frac{N^2-1}{2N(N+k)} .
\]
这个公式在第七讲会以“Wilson 线的标架相位”再次出现，那里我们用它核对三点函数的指数。

\subsection{共形块与 Ward 恒等式}

在带穿刺点 $P_1,\dots,P_n$、每个点上带表示 $R_i$ 的黎曼面上，\textbf{共形块}是满足
Ward 恒等式的“手征关联函数”的线性空间。具体说，取场 $\phi_{\lambda_i}$ 的关联函数
\[
\Psi(z_1,\dots,z_n)=\left\langle\phi_{\lambda_1}(z_1)\cdots\phi_{\lambda_n}(z_n)\right\rangle ,
\]
要求它对手征坐标全纯，并且被能-动量张量的 Ward 恒等式约束：
\[
\left\langle T(z)\prod_i\phi_i(z_i)\right\rangle
=\sum_i\left[\frac{h_i}{(z-z_i)^2}+\frac{1}{z-z_i}\partial_{z_i}\right]
\left\langle\prod_i\phi_i\right\rangle .
\]
“解这个线性系统”就是“求共形块”。对有理共形场论，解空间是有限维的；
它的维数不依赖复结构（只依赖标记点的拓扑与表示），这本身就是一个非平凡的定理。

\subsection{KZ 方程从哪里来}

把 Sugawara 表达式 (6.2) 代进 (6.5) 的左边，用 (6.1) 收缩，就得到 $\Psi$ 的一阶方程。
这是本讲最值得亲手推一遍的计算，分四步。

\textbf{第一步：收缩。}$T(z)$ 与每个 $\phi_i$ 收缩一次，留下 $J^a(z)\phi_i(z_i)\sim
t^a_i\phi_i(z_i)/(z-z_i)$；对 $\sum_a:J^aJ^a:$ 来说，这就是单极点项的来源。

\textbf{第二步：分出双极点。}把 $T$ 写成 $z$ 的展开，$\frac{1}{(z-z_i)(z-z_j)}$ 型的项
在 $z\to z_i$ 时给出 $1/(z-z_i)^2$（取 $j=i$）与 $1/(z-z_i)$（取 $j\ne i$）两类。

\textbf{第三步：双极点与共形权比较。}双极点的系数应当是 $h_i$，比较后得到
$h_i=(\lambda_i,\lambda_i+2\rho)/2(k+h^\vee)$，即 (6.4)。这一步同时验证了 Sugawara 的归一化。

\textbf{第四步：单极点给出 KZ。}单极点的系数在两边的比较给出
\begin{equation}
\kappa\,\partial_{z_i}\Psi=\sum_{j\ne i}\frac{\Omega_{ij}}{z_i-z_j}\Psi,
\qquad \kappa=k+h^\vee,\quad \Omega_{ij}=\sum_a t^a_i t^a_j . \tag{6.5}
\end{equation}
这就是 Knizhnik--Zamolodchikov（KZ）方程。系数里的 $\kappa=k+h^\vee$ 与 (6.2) 的分母是同一个；
所以第四讲中“半形式把 $k$ 平移成 $k+h^\vee$”这件事，在二维侧就是 Sugawara 公式里的那个平移——
两边说的是同一件事。

\subsection{两点与四点：解长什么样}

两点情形最简单。KZ 方程 (6.5) 只剩一项，直接积分得
\[
\Psi(z_1,z_2)=(z_1-z_2)^{-2h_j},
\]
指数里的 $k$ 通过 $h_j$ 含在 $k+h^\vee$ 里。这正是二点函数
$\langle\phi_j(z_1)\phi_j(z_2)\rangle$ 的标准形式，说明 KZ 方程的解就是共形块。

三点情形：固定 $SL(2,\CC)$ 后三个点可以放到 $0,1,\infty$，KZ 方程退化为一组代数关系；
解空间至多一维，通常记为 $N_{ijk}\in\{0,1\}$（对 $SU(2)$）。
这正是下一小节要用的“三点共形块维数=融合系数”。

四点情形开始有意思。把 $SL(2,\CC)$ 固定到 $(0,z,1,\infty)$ 后，KZ 方程变成关于 $z$ 的
\textbf{二阶 Fuchs 型方程}，解空间是二维的（一般情形）。
两个独立解对应两种“中间表示”的交换：
$s$-道把 $1,2$ 先融合、$t$-道把 $2,3$ 先融合；两套解互为解析延拓，
其单值性（monodromy）就是辫群的表示。换言之，
\textbf{KZ 方程的解空间上带着一个平坦联络，它的单值性就是编织}。
三维一侧的“纽结”从这里进来：辫群的表示给出 Jones 多项式的 $R$ 矩阵。

有一个小例子可以手算：$SU(2)$、$k=1$、四个 $j=\frac12$。
因为 $2j\le k$ 只允许 $j=0,\frac12$，而 $\frac12\otimes\frac12=0$，
所以中间表示只能是 $0$，解空间一维，四点块被唯一确定（直到整体归一化）。
这解释了为什么 $k=1$ 的 Jones 多项式特别简单：只有一个通道。

\subsection{共形块的维数与融合系数}

对有理（rational）的共形场论，共形块空间有限维，维数由融合规则给出：
三点函数空间的维数是
\[
N_{ijk}=\dim\cH_{\{P_1,P_2,P_3;R_i,R_j,R_k\}},
\]
它满足融合代数
\[
R_i\otimes R_j=\bigoplus_k N_{ij}{}^kR_k ,
\qquad N_{ij}{}^k=N_{ijk}\ \text{（用对偶度规升降指标）}.
\]
一般地，$n$ 点共形块的维数是把 $n$ 个表示按融合规则配成不变量的方案数：
\[
\dim\cH_{S^2;R_1,\dots,R_n}
=\sum_{i_1,\dots,i_{n-3}}
N_{i_1 i_2}{}^{j_1}N_{j_1 i_3}{}^{j_2}\cdots .
\]
这个“数方案”的图像在第七讲会变成“手术时插入完备态然后求和”，
两者是同一件事的两种说法。

\subsection{特征标与模不变性}

共形块在环面上求迹，得到特征标
\begin{equation}
\chi_\lambda(\tau)=\Tr_{\hat V_{k,\lambda}}\left(q^{L_0-c/24}\right),
\qquad q=e^{2\pi\ii\tau}. \tag{6.6}
\end{equation}
$q^{L_0-c/24}$ 里的 $c/24$ 是柱面上的真空能：$L_0$ 在真空上取值 $-c/24$ 而不是零，
这一点在第七讲计算标架反常时会再遇到。

模群 $\mathrm{SL}(2,\ZZ)$ 由 $\tau\to\tau+1$ 与 $\tau\to-1/\tau$ 生成，它把特征标混合起来：
\[
\chi_\lambda(\tau+1)=e^{2\pi\ii(h_\lambda-c/24)}\chi_\lambda(\tau),
\qquad
\chi_\lambda(-1/\tau)=\sum_\mu S_{\lambda\mu}\,\chi_\mu(\tau).
\]
$T$ 变换给对角相位，$S$ 变换给混合矩阵。对 $SU(2)$、level $k$，$S$ 有闭式
\begin{equation}
S_{jj'}=\sqrt{\frac{2}{k+2}}\,\sin\frac{\pi(2j+1)(2j'+1)}{k+2},
\qquad j,j'=0,\frac12,\dots,\frac k2 . \tag{6.7}
\end{equation}
模不变性（配分函数在模变换下不变）给出一个强约束：$S$ 必须是对称酉矩阵，
并且满足模群关系。用 Witten 的 $T=\mathrm{diag}\big(e^{2\pi\ii(h_j-c/24)}\big)$ 时，
\begin{equation}
S^2=(ST)^3=1 ; \tag{6.8}
\end{equation}
这一点可以手算验证。取 $k=2$，由 (6.7) 得
\[
S=\begin{pmatrix}
\frac12 & \frac{\sqrt2}{2} & \frac12\\[3pt]
\frac{\sqrt2}{2} & 0 & -\frac{\sqrt2}{2}\\[3pt]
\frac12 & -\frac{\sqrt2}{2} & \frac12
\end{pmatrix}.
\]
逐项相乘：$S^2$ 的 $(0,0)$ 元是 $\frac14+\frac12+\frac14=1$；
$(0,1)$ 元是 $\frac{\sqrt2}{4}+0-\frac{\sqrt2}{4}=0$；
$(1,1)$ 元是 $\frac12+0+\frac12=1$；
$(1,2)$ 元是 $\frac{\sqrt2}{4}+0-\frac{\sqrt2}{4}=0$。
其余同理，故 $S^2=1$。这个验算值得每个人自己做一遍：
它同时检查了 (6.7) 的归一化因子 $\sqrt{2/(k+2)}$ 与“可积表示个数 $=k+1$”两件事。

有了 $S$，融合规则就被它对角化（Verlinde）：
\begin{equation}
N_{ijk}=\sum_{l}\frac{S_{il}\,S_{jl}\,S_{kl}}{S_{0l}} . \tag{6.9}
\end{equation}
这一步的推导是第七讲的主题；在二维一侧，它的意思是“$S$ 矩阵的列是融合代数的特征向量”。
读者可以把 (6.9) 当作一句承诺：只要知道 $S$，就知道所有融合规则，
于是所有共形块维数都可以算。

\subsection{三个最小例子的数据表}

把 $k=1,2,3$ 的数据列成表，第七讲计算 $S^3$ 配分函数时会直接用到。

\begin{center}
\small
\begin{tabular}{@{}clcccc@{}}
\toprule
$k$ & 可积 $j$ & 表示个数 & $c=3k/(k+2)$ & $h_j=\dfrac{j(j+1)}{k+2}$ & $S$ 矩阵 \\
\midrule
1 & $0,\frac12$ & 2 & $1$ & $0,\ \frac14$ & $\frac{1}{\sqrt2}\begin{pmatrix}1&1\\1&-1\end{pmatrix}$ \\
2 & $0,\frac12,1$ & 3 & $\frac32$ & $0,\ \frac{3}{16},\ \frac12$ & $\frac{\sqrt2}{2}\begin{pmatrix}\frac{1}{\sqrt2}&1&\frac{1}{\sqrt2}\\1&0&-1\\ \frac{1}{\sqrt2}&-1&\frac{1}{\sqrt2}\end{pmatrix}$ \\
3 & $0,\frac12,1,\frac32$ & 4 & $\frac95$ & $0,\ \frac{3}{20},\ \frac25,\ \frac34$ & 由 (6.7) 给出 \\
\bottomrule
\end{tabular}
\end{center}

顺带给出三个检验，读者可以自己核：
\begin{enumerate}[leftmargin=*]
  \item $\sum_jS_{0j}^{\,2}=1$ 对每个 $k$ 都成立（$S$ 酉且实对称，第一行是单位向量）。
  \item $h_j$ 的公式在 $k\to\infty$ 时趋于 $j(j+1)/k$，趋于经典值；
        $c$ 在 $k\to\infty$ 时趋于 $3=\dim SU(2)$。
  \item 表中 $h_j$ 的分母是 $k+2$，不是 $k$——凡是写成 $j(j+1)/k$ 的公式都把裸 level
        与有效 level $k+h^\vee$ 混了。
\end{enumerate}

关于这些理论的“身份”，两点值得注意。
$k=1$ 时 $c=1$，$SU(2)_1$ 等价于半径 $R=\sqrt2$ 的紧致自由玻色子（自对偶点）。
$k=2$ 时 $c=3/2$，与三份 Ising 模型有相同的中心荷，但两者初级场谱不同
（$SU(2)_2$ 的共形权是 $0,\frac{3}{16},\frac12$，Ising 的是 $0,\frac{1}{16},\frac12$），
所以不是直和；把 $SU(2)_2$ 的 $U(1)$ 子代数商掉，剩下的正是 Ising 模型
（$Z_2$ 仲费米子）。$k=3$ 时 $c=9/5>1$，它不属于最小模型（最小模型的 $c<1$），
这一点常被误记。

\subsection{同一个空间}

现在回到第四讲：$\cH_\Sigma=H^0(\cM_J,L^{\otimes k})$。
Witten 的论断是：这个空间与 WZW 模型的共形块空间\textbf{是同一个向量空间}。
两者都是“带标记点曲面的共形块”，都带同一个射影平坦丛，
而且在标记点处的插入由同一个表示给出。

这个等同可以这样理解：$\cM_J$ 上 Hitchin 联络的平坦性，正好等价于共形块满足 Ward 恒等式；
射影平坦性（共形 anomaly）对应中心荷 (6.3)；而穿刺点上的抛物结构对应 WZW 的插入。
换句话说，第四讲里“$\cH_\Sigma$ 射影平坦”这句话的二维翻译，就是“存在共形块”。

\subsection{这个等同是什么级别}

需要精确地说：这是一个\textbf{对象层面}的等同——两个向量空间同构，
并且同构与射影平坦丛相容。它\textbf{不是}两个理论作为量子场论的等价：
三维理论没有边界上的动力学，二维理论有；三维理论没有引力子，
二维理论的共形块也不描述动力学。

更强的“函子性”（把曲面的缝合对应到空间的张量积）在原文里只被提到一句
“It is obvious at this point”，严格化是后来的工作（Segal 的模函子公理、
Moore--Seiberg 的缝合证明）。所以我们用“同一个向量空间”而不用“同一个理论”，
是有意的措辞。

\subsection{从 WZW 到极小模型：陪集构造}

WZW 模型是“母模型”，很多有理共形场论可以从它商出来。设 $H\subset G$ 是一个子群，
把 $G$-WZW 模型按 $H$ 的规范对称性商掉，得到陪集（coset）模型 $G/H$。
它的中心荷是
\[
c_{G/H}=c_G-c_H=\frac{k_G\dim G}{k_G+h^\vee_G}-\frac{k_H\dim H}{k_H+h^\vee_H},
\]
其中 $k_H$ 由嵌入的“水平”决定（对嵌入 $H\subset G$，水平按 Dynkin 指数的比例缩放）。

最重要的例子是 $SU(2)_k/U(1)$。它的中心荷是 $c=1-\frac{6}{(k+2)(k+3)}$，
这正是极小模型 $M(k+1,k+2)$ 的中心荷。也就是说，极小模型（包括 Ising 的 $c=\frac12$、
三态 Potts 的 $c=\frac45$）都可以写成 $SU(2)$ WZW 模型的陪集。

这个构造对本文的意义是：它说明 WZW 模型不是“一个模型”，而是一个“模型工厂”。
本文用 WZW 来解释 Chern--Simons，但 WZW 本身又生成了一大批二维模型；
这条链条（三维 CS → 二维 WZW → 二维极小模型）是 1980 年代后期共形场论的核心图景。

\subsection{与自由场的关系：玻色化}

WZW 模型不是凭空出现的，它有自由场实现。最简单的例子是 $SU(2)_1$：
它的中心荷 $c=1$，与一个自由玻色子相同。事实上，$SU(2)_1$ 的 WZW 模型等价于
一个半径为 $R=\sqrt2$ 的紧致玻色子；两个可积表示 $j=0,\frac12$ 对应玻色子的两种
“缠绕/动量”扇区，共形权分别是 $0$ 与 $\frac14$。

这个等价关系叫玻色化（bosonization）。它有两个用处。第一，它给了 WZW 模型一个
可以显式计算的模型：自由场的关联函数是已知的，所以 $SU(2)_1$ 的共形块可以手算。
第二，它解释了为什么 $SU(2)_1$ 的共形块维数与 theta 函数有关：自由玻色子的态空间
就是 theta 函数空间。

对一般的 $G$ 与 $k$，情况更复杂，但有一个一般的构造：WZW 模型可以写成
“$G$-值场 + 一个 $G$ 的 Wess--Zumino 项”，在 $k=1$ 时可以用自由场表示。
更高 level 需要引入“轨道流形”或“陪集构造”，那是另一套技术。

对本文的读者来说，记住两点就够了：第一，WZW 不是抽象的公理系统，它有具体的自由场实现；
第二，$k=1$ 与自由场的关系，是理解 $c=1$ 与 $c=k\dim G/(k+h^\vee)$ 的钥匙。

\subsection{小结}

二维一侧给出的语言是：仿射代数与 level、可积表示与截断、Sugawara 与中心荷、
共形块与 KZ 方程、融合规则与 $S$ 矩阵。它们的逻辑关系是：
\begin{itemize}[leftmargin=*]
  \item level $k$ 决定代数 $\hat{\mathfrak g}_k$，也决定可积表示只有有限个；
  \item Sugawara 构造把流代数变成 Virasoro，中心荷 $c=k\dim G/(k+h^\vee)$；
  \item 共形块是 Ward 恒等式的解，等价于 KZ 方程 (6.5) 的全纯解；
  \item 共形块维数就是融合系数 $N_{ijk}$，而 $N_{ijk}$ 由 $S$ 矩阵经 (6.9) 给出。
\end{itemize}
第七讲会把这套代数语言与三维手术对接：每一次手术对应一次融合，
每一个链环分量对应一个表示，最后得到的数就是 Jones 多项式在 $q=e^{2\pi\ii/(k+N)}$ 处的值。

\subsection*{练习}
\begin{enumerate}[leftmargin=*]
  \item 由 (6.2) 与 (6.1) 验证 $c=k\dim G/(k+h^\vee)$ 对 $SU(2)$ 给出 $3k/(k+2)$。
  \item 计算 $SU(2)$、$k=1$ 时的 $h_0,h_{1/2}$ 与 $c$，并验证 $\sum_j S_{0j}^2=1$。
  \item 证明两点 KZ 方程的解是 $(z_1-z_2)^{-2h_j}$。
  \item 验证 $SU(2)$、$k=2$ 的 $S$ 矩阵满足 $S^2=1$，并用 (6.9) 算出 $N_{1,1}{}^{0}$ 与 $N_{1,1}{}^{1}$。
  \item 说明为什么 Wess--Zumino 项不能写成 $\Sigma$ 上的二维积分；
        再说明为什么换填充改变 $S$ 恰好是 $2\pi k$ 的整数倍。
  \item 用 (6.7) 算出 $k=3$ 的 $4\times4$ $S$ 矩阵的第一行，并验证其平方和为 1。
  \item 判断：$SU(2)_2$ 与三份 Ising 模型是同一个共形场论吗？给出理由。
\end{enumerate}

\begin{quote}
\small
\textbf{本讲要点。} WZW 的 level $k$ 给出仿射代数 $\hat{\mathfrak g}_k$ 与有限多个可积表示；
Sugawara 构造给中心荷 $c=k\dim G/(k+h^\vee)$ 与共形权 (6.4)；
共形块满足 KZ 方程 (6.5)，系数 $\kappa=k+h^\vee$；
共形块维数等于融合系数，融合系数由 $S$ 矩阵经 (6.9) 给出；
三维 Hilbert 空间与共形块空间在对象层面等同。
\end{quote}
